\documentclass[12pt,a4paper]{article}

% ─── بسته‌های پایه ───────────────────────────────────────────
\usepackage{fontspec}
\usepackage{geometry}
\usepackage{xcolor}
\usepackage{colortbl}
\usepackage{booktabs}
\usepackage{longtable}
\usepackage{array}
\usepackage{tabularx}
\usepackage{multirow}
\usepackage{makecell}
\usepackage{tikz}
\usepackage{pgfplots}
\usepackage{tcolorbox}
\usepackage{fancyhdr}
\usepackage{titlesec}
\usepackage{titletoc}
\usepackage{hyperref}
\usepackage{enumitem}
\usepackage{graphicx}
\usepackage{wrapfig}
\usepackage{multicol}
\usepackage{setspace}
\usepackage{amsmath}
\usepackage{amssymb}
\usepackage{footnote}
\usepackage{etoolbox}
\usepackage{calc}
\usepackage{varwidth}
\usepackage{needspace} % FIXED: برای جلوگیری از جدا شدن عنوان از متن
\usepackage{adjustbox} % FIXED: برای تنظیم اندازه جداول و نمودارها

% ─── تنظیمات TikZ ─────────────────────────────────────────────
\usetikzlibrary{
  arrows, arrows.meta, positioning, shapes, shapes.geometric,
  shapes.multipart, shadows, decorations, decorations.pathreplacing,
  decorations.markings, patterns, backgrounds, calc, fit,
  mindmap, trees, matrix, chains, scopes,
  fadings, through
}
\pgfplotsset{compat=1.18}
\tcbuselibrary{skins, breakable}

% ─── صفحه‌آرایی ───────────────────────────────────────────────
\geometry{
  a4paper,
  top=2.8cm, bottom=2.8cm,
  right=3.5cm, left=3cm,
  headheight=1.5cm,
  headsep=0.5cm,
  footskip=1.2cm
}

% ─── پالت رنگی ────────────────────────────────────────────────
\definecolor{PrimaryDeep}{HTML}{1A237E}
\definecolor{PrimaryMid}{HTML}{283593}
\definecolor{PrimaryLight}{HTML}{3949AB}
\definecolor{AccentGold}{HTML}{F9A825}
\definecolor{AccentAmber}{HTML}{FF8F00}
\definecolor{AccentTeal}{HTML}{00695C}
\definecolor{AccentRed}{HTML}{B71C1C}
\definecolor{AccentGreen}{HTML}{1B5E20}
\definecolor{AccentPurple}{HTML}{4A148C}
\definecolor{NeutralLight}{HTML}{ECEFF1}
\definecolor{NeutralMid}{HTML}{B0BEC5}
\definecolor{NeutralDark}{HTML}{455A64}
\definecolor{TextMain}{HTML}{212121}
\definecolor{TextSecond}{HTML}{424242}
\definecolor{BgWarm}{HTML}{FFF8E1}
\definecolor{BgCool}{HTML}{E8EAF6}
\definecolor{BgTeal}{HTML}{E0F2F1}

\definecolor{EraAncient}{HTML}{4E342E}
\definecolor{EraMedieval}{HTML}{37474F}
\definecolor{EraEarlyMod}{HTML}{1565C0}
\definecolor{EraModern}{HTML}{2E7D32}
\definecolor{EraContemp}{HTML}{6A1B9A}
\definecolor{EraPostmod}{HTML}{AD1457}

% ─── جعبه‌های محتوایی ─────────────────────────────────────────
\tcbset{
  mybox/.style={
    enhanced, breakable,
    colback=BgCool, colframe=PrimaryMid,
    fonttitle=\bfseries\large,
    attach boxed title to top right={yshift=-2mm, xshift=-3mm},
    boxed title style={colback=PrimaryMid, colframe=PrimaryDeep},
    top=4mm, bottom=4mm,
    before skip=12pt, after skip=12pt,
  },
  defbox/.style={
    enhanced, breakable,
    colback=BgWarm, colframe=AccentGold,
    fonttitle=\bfseries,
    left=4mm, right=4mm,
    borderline west={3pt}{0pt}{AccentAmber},
    sharp corners,
    before skip=12pt, after skip=12pt,
  },
  wavebox/.style={
    enhanced, breakable,
    colback=BgTeal, colframe=AccentTeal,
    fonttitle=\bfseries,
    rounded corners,
    drop shadow,
    before skip=12pt, after skip=12pt,
  },
  quotebox/.style={
    enhanced, breakable,
    colback=NeutralLight, colframe=NeutralDark,
    fontupper=\itshape,
    left=8mm,
    borderline west={4pt}{0pt}{NeutralMid},
    sharp corners=south,
    before skip=12pt, after skip=12pt,
  },
  enemybox/.style={
    enhanced, breakable,
    colback=red!5, colframe=AccentRed,
    fonttitle=\bfseries,
    borderline west={3pt}{0pt}{AccentRed},
    before skip=12pt, after skip=12pt,
  },
  warnbox/.style={
    enhanced, breakable,
    colback=orange!5, colframe=AccentAmber,
    fonttitle=\bfseries,
    borderline west={3pt}{0pt}{AccentAmber},
    before skip=12pt, after skip=12pt,
  },
  successbox/.style={
    enhanced, breakable,
    colback=green!5, colframe=AccentGreen,
    fonttitle=\bfseries,
    borderline west={3pt}{0pt}{AccentGreen},
    before skip=12pt, after skip=12pt,
  },
}

% ─── سربرگ و پابرگ ────────────────────────────────────────────
\pagestyle{fancy}
\fancyhf{}
\renewcommand{\headrulewidth}{0.6pt}
\renewcommand{\footrulewidth}{0.3pt}
\fancyhead[R]{\small\color{PrimaryMid}\rightmark}
\fancyhead[L]{\small\color{NeutralDark}اخلاق براندازی}
\fancyfoot[C]{\small\thepage}

% ─── عنوان‌بندی بخش‌ها ────────────────────────────────────────
\titleformat{\section}[block]
  {\Large\bfseries\color{PrimaryDeep}}
  {\thesection}{1em}{}
  [{\color{AccentGold}\titlerule[1.5pt]}]

\titleformat{\subsection}[block]
  {\large\bfseries\color{PrimaryMid}}
  {\thesubsection}{0.8em}{}
  [{\color{NeutralMid}\titlerule[0.6pt]}]

\titleformat{\subsubsection}[runin]
  {\normalsize\bfseries\color{AccentTeal}}
  {\thesubsubsection}{0.6em}{}[\ —\ ]

% ─── فاصله‌گذاری ──────────────────────────────────────────────
\titlespacing*{\section}{0pt}{18pt}{10pt}
\titlespacing*{\subsection}{0pt}{14pt}{6pt}
\setlength{\parindent}{1em}
\setlength{\parskip}{0.4em}
\onehalfspacing

% ─── hyperref ─────────────────────────────────────────────────
\hypersetup{
  colorlinks  = true,
  linkcolor   = PrimaryMid,
  citecolor   = AccentTeal,
  urlcolor    = AccentAmber,
  pdftitle    = {Ethics of Subversion},
  pdfauthor   = {},
  pdfsubject  = {Political Ethics and Philosophy of Resistance},
  bookmarks   = true,
  pdfpagemode = UseOutlines,
  pdfencoding = unicode,
}

\usepackage{xepersian} % FIXED: آخرین بسته

% ─── فونت‌ها ──────────────────────────────────────────────────
\settextfont[Scale=1.05]{Vazirmatn}
\setlatintextfont[Scale=0.95]{Linux Libertine}
\setdigitfont[Scale=0.95]{Vazirmatn}

% ─── دستورات کمکی ─────────────────────────────────────────────
\newcommand{\concept}[1]{\textbf{\color{PrimaryMid}#1}}
\newcommand{\thinker}[1]{\textit{\color{AccentTeal}#1}}
\newcommand{\era}[1]{\textbf{\color{AccentAmber}#1}}
\newcommand{\enemy}[1]{\textbf{\color{AccentRed}#1}}
\newcommand{\lat}[1]{\lr{#1}}
\newcommand{\src}[1]{\textsuperscript{\scriptsize\color{NeutralDark}[#1]}}

% ──────────────────────────────────────────────────────────────
\begin{document}

% ========== صفحه عنوان ==========
\begin{titlepage}
\begin{center}
\vspace*{2cm}
{\color{PrimaryDeep}\rule{\linewidth}{2pt}}
\vspace{1cm}

{\Huge\bfseries\color{PrimaryDeep} اخلاق براندازی}

\vspace{0.5cm}

{\Large\color{PrimaryMid} مبانی نظری، ابعاد عملی و چالش‌های اخلاقی}

\vspace{0.3cm}

{\Large\color{AccentTeal} مقاومت و مبارزه با استبداد}

\vspace{0.5cm}

{\large\color{NeutralDark} با نگاهی ویژه به تجربهٔ ایران}

\vspace{1cm}
{\color{AccentGold}\rule{0.5\linewidth}{1pt}}
\vspace{2cm}

{\large\color{TextSecond}
حاشیه علیه متن\\[0.3cm]
حق عصیان، اخلاق دشمنی و راه میانهٔ آزادی‌خواهی
}

\vspace{2cm}

{\normalsize\color{NeutralDark} \today}

\vspace{1cm}
{\color{PrimaryDeep}\rule{\linewidth}{2pt}}
\end{center}
\end{titlepage}

% ========== چکیده ==========
\thispagestyle{empty}
\begin{tcolorbox}[mybox, title=چکیده]
این مقاله به بررسی جامع \concept{اخلاق براندازی} و مبارزه با نظام‌های استبدادی 
می‌پردازد. پرسش محوری آن است که گروه‌ها و ملت‌هایی که تحت ستم سیستماتیک یک رژیم 
سیاسی قرار دارند، از چه حقوقی برای براندازی برخوردارند و در اِعمال این حق، چه 
اصول و ملاحظات اخلاقی‌ای را باید رعایت کنند — به‌گونه‌ای که نه به بی‌عملی و سستی 
منجر شود و نه به تبدیل‌شدن مبارز به همان شرّی که علیه آن مبارزه می‌کند.

مقاله ابتدا در یک نمای شماتیک، چارچوب مفهومی بحث را ترسیم می‌کند و سپس به بررسی 
تفصیلی مبانی نظری، دیدگاه‌های مکاتب فلسفی و ایدئولوژیک مختلف، اخلاق جنگ و دشمنی، 
مقاومت خشونت‌پرهیز، و مسئلهٔ حیاتیِ تناسب و محدودیت در کاربرد قدرت می‌پردازد. 
مورد ایران به‌عنوان نمونه‌ای زنده از چالش‌های اخلاقی مبارزه با استبداد بررسی می‌شود. 
در پایان، پیوست‌هایی شامل مرور تاریخی، موج‌شناسی جنبش‌های مقاومت و کتاب‌شناسی 
جامع ارائه شده است.

\textbf{واژگان کلیدی:} اخلاق براندازی، حق عصیان، مقاومت مدنی، اخلاق جنگ، 
اخلاق دشمنی، مبارزهٔ خشونت‌پرهیز، دفاع مشروع، نظریهٔ جنگ عادلانه، ایران.
\end{tcolorbox}

\newpage
\tableofcontents
\newpage

% ╔══════════════════════════════════════════════════════════════╗
% ║               بخش صفر: نمای کلی شماتیک                    ║
% ╚══════════════════════════════════════════════════════════════╝

\section*{بخش صفر: نمای کلی شماتیک}
\addcontentsline{toc}{section}{بخش صفر: نمای کلی شماتیک}

\subsection*{نقشهٔ مفهومی اخلاق براندازی}

این نقشه، ساختار کلی مسائل اخلاقی مرتبط با براندازی را نشان می‌دهد. هر شاخه
در بخش‌های بعدی مقاله به تفصیل بررسی خواهد شد.

\begin{figure}[htbp]
\centering
\begin{adjustbox}{max width=\linewidth}
\begin{tikzpicture}[
    mindmap,
    every node/.style={concept, circular drop shadow, minimum size=0pt,
      text width=2.2cm, font=\small\bfseries, align=center},
    root concept/.append style={
      concept color=PrimaryDeep, fill=PrimaryDeep, text=white,
      font=\large\bfseries, text width=3cm, minimum size=3cm},
    level 1 concept/.append style={
      font=\small\bfseries, text width=2cm, minimum size=1.8cm,
      sibling angle=60},
    level 2 concept/.append style={
      font=\tiny\bfseries, text width=1.5cm, minimum size=1.2cm,
      sibling angle=45},
    grow cyclic,
  ]
  \node[root concept] {\rl{اخلاق}\\\rl{براندازی}}
    child[concept color=AccentTeal] { node {\rl{حقوق}\\\rl{بنیادین}}
      child { node {\rl{حق عصیان}} }
      child { node {\rl{دفاع مشروع}} }
      child { node {\rl{حق آزادی}} }
    }
    child[concept color=AccentGold!80!black] { node {\rl{چارچوب}\\\rl{اخلاقی}}
      child { node {\rl{وظیفه‌گرایی}} }
      child { node {\rl{فایده‌گرایی}} }
      child { node {\rl{فضیلت‌محوری}} }
    }
    child[concept color=AccentRed!80!black] { node {\rl{روش‌های}\\\rl{مبارزه}}
      child { node {\rl{خشونت‌پرهیز}} }
      child { node {\rl{مسلحانه}} }
      child { node {\rl{ترکیبی}} }
    }
    child[concept color=AccentPurple] { node {\rl{قیدها و}\\\rl{محدودیت‌ها}}
      child { node {\rl{تناسب}} }
      child { node {\rl{تمایز}} }
      child { node {\rl{ضرورت}} }
    }
    child[concept color=AccentGreen] { node {\rl{پیامدها و}\\\rl{خطرها}}
      child { node {\rl{ابتذال شر}} }
      child { node {\rl{بی‌عملی}} }
      child { node {\rl{عدالت نوین}} }
    }
    child[concept color=EraEarlyMod] { node {\rl{مورد}\\\rl{ایران}}
      child { node {\rl{سرکوب}} }
      child { node {\rl{جنبش‌ها}} }
      child { node {\rl{چشم‌انداز}} }
    };
\end{tikzpicture}
\end{adjustbox}
\caption*{نقشهٔ مفهومی اخلاق براندازی}
\end{figure}

\subsection*{مغز مطلب در یک نگاه}

\begin{tcolorbox}[defbox, title={\textbf{پرسش بنیادین}}]
وقتی گروهی تحت ستم سیستماتیک یک رژیم سیاسی هستند:
\begin{enumerate}[label=\textbf{\arabic*.}]
\item \textbf{آیا حق براندازی دارند؟} — بله. تقریباً تمامی سنت‌های فکری بزرگ، 
از الاهیات مسیحی و اسلامی گرفته تا فلسفهٔ لیبرال و مارکسیسم، در شرایط خاصی 
حق مقاومت و عصیان را به رسمیت می‌شناسند.
\item \textbf{این حق مطلق است؟} — خیر. هر مجوز نظری دوطرفه است: اگر خشونت 
را بی‌قید و شرط مجاز بدانیم، همان مجوز در دست طرف مقابل هم قرار می‌گیرد.
\item \textbf{چه محدودیت‌های اخلاقی‌ای وجود دارد؟} — محدودیت‌هایی که از 
\concept{تناسب}، \concept{ضرورت}، \concept{تمایز} (میان مبارز و غیرمبارز)، 
و \concept{کرامت انسانی} سرچشمه می‌گیرند.
\item \textbf{آیا این محدودیت‌ها به بی‌عملی نمی‌انجامند؟} — لزوماً نه. کلید، 
یافتن \concept{راه میانه} میان ساده‌سازی افراطی (هر کاری مجاز است) و 
پیچیده‌سازی فلج‌کننده (هیچ کاری مجاز نیست) است.
\end{enumerate}
\end{tcolorbox}

\subsection*{جدول خلاصهٔ چارچوب تحلیلی}

\begin{table}[htbp]
\centering
\small
\begin{tabularx}{\linewidth}{|>{\bfseries\raggedleft}p{2.5cm}|X|X|X|}
\hline
\rowcolor{BgCool}
\textbf{بُعد} & \textbf{پرسش محوری} & \textbf{خطر ساده‌سازی} & \textbf{خطر پیچیده‌سازی} \\
\hline
حق عصیان & آیا مردم حق سرنگونی حاکم ظالم را دارند؟ & هر نارضایتی = حق انقلاب & هیچ شرایطی کافی نیست \\
\hline
روش مبارزه & چه ابزارهایی مجاز است؟ & همه‌چیز مجاز است & فقط اعتراض مؤدبانه \\
\hline
تناسب & شدت عمل چقدر باشد؟ & تلافی نامحدود & واکنش صفر \\
\hline
کرامت انسانی & حرمت دشمن چه جایگاهی دارد؟ & دشمن انسان نیست & هرگز نباید آسیبی رسید \\
\hline
پیامدها & نتیجهٔ مطلوب چیست؟ & فقط پیروزی مهم است & فقط فرایند مهم است \\
\hline
\end{tabularx}
\caption*{چارچوب تحلیلی: میان دو پرتگاه}
\end{table}

\subsection*{قاعدهٔ طلایی: دوطرفه‌بودن هر مجوز نظری}

\begin{tcolorbox}[enemybox, title={\textbf{هشدار بنیادین}}]
\textbf{هر مجوز نظری‌ای دوطرفه است.} اگر بگوییم «خشونت علیه ظالم مجاز است» 
بدون تعریف دقیق «ظالم» و «خشونت مجاز»، همان استدلال را حکومت مستبد هم علیه 
معترضان به کار خواهد بست. مگر آنکه محدودیت‌ها دقیق، موجه و مقبول باشد. 
این بزرگ‌ترین چالش اخلاق براندازی است: چگونه می‌توان مرزی ترسیم کرد که 
مبارزان آزادی‌خواه را توانمند سازد اما از سوءاستفاده و تبدیل‌شدن به همان 
شرّی که علیه‌اش می‌جنگند، جلوگیری کند؟
\end{tcolorbox}

\subsection*{مسیر مقاله}

\begin{tcolorbox}[wavebox, title={\textbf{ساختار کلی مقاله}}]
\begin{enumerate}[label=\textbf{بخش \arabic*:}]
\item \textbf{مبانی نظری} — تعریف مفاهیم، تمایز 
\lat{Ethics} و \lat{Morality}، ریشه‌های تاریخی حق عصیان
\item \textbf{دیدگاه‌های فلسفی و ایدئولوژیک} — از ارسطو تا فانون، از لاک تا 
شریعتی، از گاندی تا شارپ
\item \textbf{اخلاق کاربردی براندازی} — اخلاق دشمنی، اخلاق جنگ، تناسب، 
خویشتنداری، تفاوت خشونت دولتی و مردمی
\item \textbf{مورد ایران} — تحلیل اخلاقی جنبش‌ها و اعتراضات، دوگانه‌ها و 
چالش‌ها
\item \textbf{پیچیدگی‌ها و دام‌ها} — ساده‌سازی و پیچیده‌سازی، اخلاقِ مانع و 
اخلاقِ یاری‌رسان، راه میانه
\item \textbf{نتیجه‌گیری}
\end{enumerate}
\textbf{پیوست‌ها:} مرور تاریخی و موج‌شناسی — جدول تطبیقی جنبش‌ها — کتاب‌شناسی جامع
\end{tcolorbox}


% ╔══════════════════════════════════════════════════════════════╗
% ║        بخش اول: مبانی نظری                                ║
% ╚══════════════════════════════════════════════════════════════╝

\newpage
\section{مبانی نظری: مفاهیم، تمایزها و ریشه‌ها}

\subsection{تعریف مفاهیم کلیدی}

\subsubsection{براندازی، انقلاب، مقاومت}
پیش از هر بحثی، تعریف دقیق واژگان ضروری است. «براندازی» 
(\lat{Subversion/Overthrow}) به معنای تلاش سازمان‌یافته برای سرنگونی نظم 
سیاسی موجود و جایگزینی آن با نظمی نوین است. این مفهوم با «انقلاب» 
(\lat{Revolution}) همپوشانی دارد اما لزوماً هم‌معنا نیست: انقلاب غالباً بار 
دگرگونی بنیادین اجتماعی-اقتصادی هم دارد، در حالی که براندازی می‌تواند صرفاً 
تغییر نظام حکومتی را هدف بگیرد. «مقاومت» (\lat{Resistance}) مفهومی عام‌تر 
است که طیفی از نافرمانی مدنی تا مبارزهٔ مسلحانه را دربرمی‌گیرد.

«نافرمانی مدنی» (\lat{Civil Disobedience}) — آن‌گونه که \thinker{جان رالز} 
تعریف می‌کند — عملی عمومی، خشونت‌پرهیز و وجدانی است که با هدف تغییر قانون یا 
سیاست حکومت و با پذیرش مجازات قانونی صورت می‌گیرد.\footnote{\lr{Rawls, John. 
\textit{A Theory of Justice}. 1971, §55.}} اما وقتی حکومت اساساً مشروعیت 
ندارد و فضایی برای اعتراض قانونی نگذاشته، آیا نافرمانی مدنی کافی است؟

«دفاع مشروع» (\lat{Legitimate Self-Defense}) در حقوق بین‌الملل و اکثر 
نظام‌های حقوقی به رسمیت شناخته شده و حق استفاده از زور متناسب برای دفع تجاوز 
را می‌دهد. پرسش آن است که آیا مردمی که حکومتشان خودْ متجاوز به حقوق آنهاست، 
می‌توانند از این مفهوم بهره ببرند.

\subsubsection{ستم سیستماتیک}
ستم سیستماتیک (\lat{Systematic Oppression}) با ستم موردی و اتفاقی متفاوت است. 
در ستم سیستماتیک، نهادها، قوانین و ساختارهای قدرت به‌شکلی هماهنگ و مداوم علیه 
گروهی عمل می‌کنند. این تمایز اهمیت اخلاقی دارد: ستم سیستماتیک پاسخ‌های 
سیستماتیک‌تری را اخلاقاً توجیه‌پذیرتر می‌سازد.

\subsection{تمایز اخلاق و اخلاقیات: \texorpdfstring{\lat{Ethics}}{Ethics} و 
\texorpdfstring{\lat{Morality}}{Morality}}

\begin{tcolorbox}[defbox, title={\textbf{تمایزی کلیدی}}]
در ادبیات فلسفی، تمایزی ظریف اما مهم میان دو مفهوم وجود دارد:
\begin{itemize}[rightmargin=1cm]
\item \concept{اخلاقیات} (\lat{Morality}): مجموعه‌ای از هنجارها، ارزش‌ها و 
قواعد رفتاری که در یک جامعه یا سنت پذیرفته شده‌اند. اخلاقیات غالباً پیشینی، 
سنتی و مبتنی بر اجماع هستند.
\item \concept{اخلاق} (\lat{Ethics}): تأمل فلسفی و نقادانه دربارهٔ درستی و 
نادرستی، خوبی و بدی. اخلاق می‌تواند خودِ اخلاقیاتِ رایج را به چالش بکشد.
\end{itemize}
\end{tcolorbox}

این تمایز در بحث براندازی اهمیتی ویژه دارد. چه بسا \textit{اخلاقیات} رایج 
(مثلاً «اطاعت از حاکم»، «پرهیز از فتنه») دقیقاً همان ابزاری باشند که حکومت 
مستبد برای تداوم سلطه‌اش به کار می‌بندد. \textit{اخلاقِ} فلسفی اما می‌تواند 
از موضع نقادانه بپرسد: آیا این هنجارها واقعاً عادلانه‌اند یا صرفاً در خدمت 
قدرت؟ آیا «آرامش» و «نظم» ارزشمندند حتی وقتی بر ستم استوارند؟

\thinker{مارتین لوتر کینگ} در \textit{نامه از زندان بیرمنگام} 
(\lat{Letter from Birmingham Jail}، ۱۹۶۳) نوشت: «صلح واقعی تنها غیاب تنش

ادامهٔ مقاله: